% SPDX-License-Identifier: CC-BY-4.0
\documentclass[11pt,letterpaper]{article}

\usepackage[T1]{fontenc}
\usepackage{amsmath,amssymb}
\usepackage{tikz}
\usepackage{pgfplots}
\pgfplotsset{compat=1.18}
\usepackage{algorithm}
\usepackage{algpseudocode}
\usepackage{booktabs}
\usepackage{enumitem}
\usepackage[hidelinks]{hyperref}

\newcommand{\N}{\mathbb{N}}
\newcommand{\Z}{\mathbb{Z}}
\newcommand{\Q}{\mathbb{Q}}
\newcommand{\R}{\mathbb{R}}
\newcommand{\C}{\mathbb{C}}
\newcommand{\Map}[3]{#1 \colon #2 \to #3}
\newcommand{\Mod}[1]{#1\text{-}\mathrm{Mod}}
\newcommand{\GroupHom}{\mathrm{Hom}_{\mathbf{Grp}}}
\newcommand{\RingHom}{\mathrm{Hom}_{\mathbf{Ring}}}

\newtheorem{theorem}{Theorem}
\newtheorem{lemma}[theorem]{Lemma}
\newtheorem{definition}[theorem]{Definition}
\newtheorem{remark}[theorem]{Remark}
\newenvironment{proof}{\noindent\textit{Proof.}}{\hfill\(\square\)}

\title{Guarded additive cycles}
\author{}
\date{\today}

\begin{document}
\maketitle

\begin{abstract}
We classify every positive cycle of the additive raw map by a filtered
permutation of residues modulo the divisor-coprime part of its offset.
The filter is an exact integer window, including when the extracted scaling
factor is not coprime to the divisor. Both reconstruction directions preserve
the original instruction count and least period. We derive the grouped raw
period formula, identify the precise multiplicative order, and give a finite
cycle enumeration with separate preprocessing, arithmetic, and output costs.
\end{abstract}

\section{Raw steps and first-division endpoints}

Fix integers $c\ge1$ and $D\ge2$. On the positive integers, let
\[
T(n)=\begin{cases}n/D,&D\mid n,\\n+c,&D\nmid n.\end{cases}
\]
Division has precedence. The iterate $T^r$ counts $r$ original instructions.
Put
\[
g=\gcd(c,D),\quad L=D/g,\quad B=c(L-1),\quad
N=\left\lfloor\frac{B}{D-1}\right\rfloor.
\]
We use the access and core results of
\href{../../QCD.Additive.001/tex/Additive.tex}{QCD.Additive.001}.
A state $n$ is accessible precisely when $g\mid n$. In that case there is a
unique index $0\le\ell<L$ with $D\mid n+\ell c$, every earlier state is
nondivisible, and
\[
H(n)=\frac{n+\ell c}{D}=T^{\ell+1}(n).
\]
An H edge contains exactly one division, even if another division immediately
follows. An inaccessible state has the exact raw tail $n+jc$.
Every positive seed reaches a state at most $N$ or an inaccessible state;
accessible endpoints within $1,\ldots,N$ stay there. Also $N\le c$.

For a path $x_0,\ldots,x_h$ of H endpoints, write its indices as
$\ell_0,\ldots,\ell_{h-1}$ and its raw clock as
\[
R_h=\sum_{i=0}^{h-1}(\ell_i+1).
\]
Its exact equations are $Dx_{i+1}=x_i+\ell_i c$.
An H cycle is primitive if $x_h=x_0$ and $x_i\ne x_0$ for $0<i<h$.
A raw cycle has the analogous least positive return period. Cycles are
identified up to cyclic change of starting phase, not by their chosen clocks.

\section{Offset decomposition and the exact window}

% Gcd removal is a terminating construction, not a factorization assumption.
\begin{theorem}[Divisor-supported factor extraction]
\label{thm:decomposition}
There are positive integers $s,t$ and an integer $k\ge0$ such that
\[
c=st,\qquad \gcd(t,D)=1,\qquad s\mid D^k.
\]
Every divisor of $c$ coprime to $D$ divides $t$. In particular $t$ is the
unique divisibility-maximal divisor of $c$ coprime to $D$.
\end{theorem}
\begin{proof}
Apply strong induction to the positive offset. If $\gcd(c,D)=1$, take
$s=1,t=c,k=0$. Otherwise put $a=\gcd(c,D)>1$. The positive quotient $c/a$
is smaller than $c$, so write $c/a=s't$, where $t$ is coprime to $D$ and
$s'\mid D^{k'}$. Then $s=as'$ satisfies $st=c$ and $s\mid D^{k'+1}$,
since $a\mid D$. This proves existence by a strictly decreasing induction.
If $u\mid c$ and $\gcd(u,D)=1$, then $\gcd(u,s)=1$, because $s\mid D^k$.
Euclid's lemma applied to $u\mid st$ gives $u\mid t$.
If another factor has the same maximality property, each divides the other;
positive divisibility in both directions gives equality.
\end{proof}

Define the normalized window
\[
M=\left\lfloor\frac{t(L-1)}{D-1}\right\rfloor.
\]
Since $L\le D$, we have $M\le t$. There can be no positive representative
when $M=0$. No condition $\gcd(s,D)=1$ is imposed.

% The supported factor is recovered from a return equation, not cancelled through D.
\begin{theorem}[Necessary cycle window]
\label{thm:necessary}
Every positive closed H path of length $h>0$ satisfies
\[
s\mid x_i,\qquad x_i\le N\qquad(0\le i\le h).
\]
Writing $x_i=sm_i$ therefore gives $1\le m_i\le M$ and
\[
Dm_{i+1}\equiv m_i\pmod t\qquad(0\le i<h).
\]
Conversely the scaled bound satisfies $sm\le N$ if and only if $m\le M$.
\end{theorem}
\begin{proof}
Successively multiplying the edge equations gives
\[
D^h x_h=x_0+cK_h,\qquad K_h=\sum_{i=0}^{h-1}D^i\ell_i.
\]
Closure yields $(D^h-1)x_0=cK_h$. Since $s\mid c$ and
$\gcd(s,D^h-1)=1$, we obtain $s\mid x_0$. The coprimality follows from
$s\mid D^k$ and $\gcd(D,D^h-1)=1$.
Starting with $s\mid x_h=x_0$, the equation
$x_i=Dx_{i+1}-\ell_i c$ propagates divisibility backwards to every vertex.
This subtracts multiples of $s$; it does not cancel $D$.

Let $\delta(x)=(D-1)x-B$. The core estimate gives
$D\delta(x_{i+1})\le\delta(x_i)$. Iterating and using closure gives
$D^h\delta(x_0)\le\delta(x_0)$. Since $D^h>1$, this forces
$\delta(x_0)\le0$ and hence $x_0\le N$. Core invariance bounds all subsequent
vertices. For a nonnegative integer $m$,
\[
\begin{aligned}
sm\le N &\iff (D-1)sm\le st(L-1)\\
 &\iff (D-1)m\le t(L-1)\iff m\le M.
\end{aligned}
\]
Here $s>0$ justifies cancellation. Dividing the edge equation by $s$ gives
$Dm_{i+1}=m_i+\ell_i t$, which proves the residue relation.
\end{proof}

% The window makes a residue edge a first-division edge with every early guard.
\begin{theorem}[Legal modular edge]
\label{thm:edge}
Suppose $1\le m,m'\le M$ and $Dm'\equiv m\pmod t$. Then
\[
\ell=\frac{Dm'-m}{t}
\]
is a nonnegative integer less than $L$. The states $sm$ and $sm'$ are joined
by exactly $\ell$ additions and one division. No earlier division is legal.
\end{theorem}
\begin{proof}
If $Dm'<m$, the positive difference $m-Dm'$ is a multiple of $t$, but it is
strictly smaller than $t$ because $m\le M\le t$ and $Dm'\ge1$. This is
impossible. Thus $Dm'-m\ge0$, and the congruence gives its divisibility by $t$.
Furthermore
\[
\ell t=(D-1)m'+m'-m
\le t(L-1)+t-1<Lt.
\]
The first inequality uses $(D-1)m'\le t(L-1)$, $m'\le t$, and $m\ge1$.
Consequently $0\le\ell<L$. Scaling gives
$Dsm'=sm+\ell c$. Thus $g\mid sm$, and access-index uniqueness identifies
$\ell$ with its first division index. Every smaller index fails the division
guard. The endpoint is $sm'$ at raw time $\ell+1$.
\end{proof}

\section{Raw recovery and primitive cycles}

% Segmentation preserves the same raw path and its original total clock.
\begin{theorem}[Exact segmentation]
\label{thm:segmentation}
Every positive raw return of period $p>0$ has a forward phase $x\le N$.
Its length-$p$ raw segment from that phase decomposes into a closed H path
with exactly the same total raw clock $R_h=p$.
A raw cycle's least return period is unchanged by a forward phase shift.
\end{theorem}
\begin{proof}
The core-or-trap reduction supplies a reached core state or a trap. A reached
trap has a strictly increasing tail $n+jc$ and cannot return, whereas every
forward phase of a period-$p$ orbit still returns after $p$ steps. Therefore
a core phase exists and is positive.

We first segment any raw path starting in the positive core and ending there,
by strong induction on its raw length $r$. Length zero gives the empty path.
For $r>0$, the first $r$ instructions cannot all be additions: their endpoint
would be $x+rc>c$, contrary to its bound $N\le c$. Let $\ell<r$ be the first
division index. By access uniqueness, $\ell<L$, and the first endpoint is
$H(x)$. It is positive and in the core. The remaining segment has length
$r-(\ell+1)<r$, so the induction hypothesis segments it. Prepending this
first legal block recovers the same original segment and clock $r$.
Applying this to a return gives a closed H path; its length is positive
because its raw clock is $p>0$.

For phase transport, let $x=T^u(n)$ and $T^p(n)=n$. Commutation of powers
gives $T^p(x)=x$. If $T^v(x)=x$, choose $a$ with $ap\ge u$ and apply
$T^{ap-u}$ to this equality. Since $T^{ap}(n)=n$, it follows that
$T^v(n)=n$. Thus the sets of return times at these two phases coincide.
This argument uses a complementary iterate along this particular periodic
orbit, not global injectivity of $T$.
\end{proof}

% The two maps preserve least periods and can be inverted using division endpoints.
\begin{theorem}[Cycle classification with primitive recovery]
\label{thm:classification}
The positive primitive raw cycles, up to cyclic phase, correspond bijectively
to the primitive cycles of multiplication by $D^{-1}$ modulo $t$ whose
positive representatives all lie in $1,\ldots,M$.
Residue zero is represented by $t$, not by zero. Each modular edge expands
as in Theorem~\ref{thm:edge}; every raw cycle is recovered in this way, with
its original least raw period.
\end{theorem}
\begin{proof}
Coprimality makes multiplication by $D^{-1}$ a permutation of $\Z/t\Z$.
Positive representatives $1,\ldots,t$ are unique, including the zero class.
Theorem~\ref{thm:edge} expands every edge of a modular cycle wholly within
the window into a legal raw block. Its concatenation is a raw return with
clock $R_h$.

Suppose the modular cycle is primitive but a raw return $0<r<R_h$ exists.
Its endpoint is the same positive core state, so
Theorem~\ref{thm:segmentation} segments it into $k$ H blocks of total clock
$r$. The bounded first access index is unique. Induction therefore identifies
the common prefix of this path with that of the original H path, including
both states and clocks. If $k<h$, its endpoint is an earlier modular return,
contrary to primitiveness. If $k\ge h$, positivity of every block length gives
$r\ge R_h$, also impossible. Thus the expanded raw period is least.

Conversely, take a primitive raw cycle and shift to a core phase.
Segmentation and Theorem~\ref{thm:necessary} produce a closed modular path
inside the window with raw clock equal to its least period. An earlier
modular return would reconstruct a smaller positive raw return, so the
modular path is primitive. Phase transport preserves the original least
period.

Finally, each division endpoint lies at most $N\le c$. Every intermediate
addition state is $x+jc>c$ for $j\ge1$. Thus the raw cycle's states at most
$c$ are exactly its H endpoints. Dividing these by $s$ recovers its modular
cycle, up to phase. Conversely the modular edges and their unique indices
fix the full expansion. Hence the correspondence is both injective and
surjective on cyclic phases.
\end{proof}

\section{Grouped clocks and marked multiplicative orders}

% Summing before dividing is essential; individual rounded contributions are invalid.
\begin{theorem}[Exact grouped raw period]
\label{thm:clock}
For a closed H path of length $h$,
\[
c\sum_{i=0}^{h-1}\ell_i=(D-1)\sum_{i=0}^{h-1}x_i.
\]
Its raw clock, and for a primitive path its least raw period, is
\[
P=h+\frac{(D-1)\sum_{i=0}^{h-1}x_i}{c}
 =h+\frac{(D-1)\sum_{i=0}^{h-1}m_i}{t}.
\]
Both quotients are integers. The sums must be grouped before division.
\end{theorem}
\begin{proof}
Sum $Dx_{i+1}=x_i+\ell_i c$. Closure identifies
$\sum_{i<h}x_{i+1}=\sum_{i<h}x_i$, so subtraction yields the first identity.
Because $c>0$, its quotient by $c$ is exactly $\sum\ell_i$. Substitute this
in $P=\sum(\ell_i+1)$. For $x_i=sm_i$ and $c=st$, cancellation of the
positive factor $s$ gives $t\sum\ell_i=(D-1)\sum m_i$, establishing both
integrality and the normalized formula. Primitiveness follows from
Theorem~\ref{thm:classification}.
\end{proof}

Our multiplicative-order convention is the order in the finite unit group;
standard group terminology is as in Dummit--Foote~\cite{DummitFoote04}.
The next proof also gives the finite bound directly.

% The effective modulus depends on the marked state, not solely on the ambient offset.
\begin{theorem}[Exact order and finite length bound]
\label{thm:order}
Let $m_0,\ldots,m_h=m_0$ be a primitive modular cycle. Put
\[
a=\gcd(m_0,t),\qquad \lambda=t/a.
\]
Then $h=\operatorname{ord}_{\lambda}(D)$ and $1\le h\le\lambda\le t\le c$.
For $\lambda=1$, the order is one.
\end{theorem}
\begin{proof}
Iterating the residue equations gives $D^i m_i\equiv m_0\pmod t$.
Since $D$ is a unit and positive representatives up to $t$ are unique,
\[
m_i=m_0
\iff D^i m_0\equiv m_0\pmod t
\iff D^i\equiv1\pmod\lambda.
\]
For the last equivalence write $m_0=av$ and $t=a\lambda$ with
$\gcd(v,\lambda)=1$, cancel $a$, and use Euclid's lemma to cancel $v$.
The least positive return is therefore exactly the asserted order.
If two powers $D^i,D^j$, with $0\le i<j<h$, were equal modulo $\lambda$,
cancelling the unit $D^i$ would give $D^{j-i}\equiv1$, contradicting
minimality. There are $h$ distinct residues among the $\lambda$ residue
classes, so $h\le\lambda$. Finally $\lambda\le t$ and $t\le st=c$.
At modulus one the unit group is a singleton and the least positive order
is one, consistent with the representative $t$ of its zero class.
\end{proof}

\section{Finite enumeration and costs}

% Exhaustion scans full permutation orbits and rejects any orbit failing the window.
\begin{theorem}[Finite complete cycle enumeration]
\label{thm:enumeration}
After factor extraction and modular inversion, all positive raw cycles can
be enumerated by at most $t$ modular-state visits. One must keep a permutation
orbit if and only if every positive representative lies at most $M$.
The total number of surviving representatives is at most $t$, and their
complete raw expansions have total length at most $tL$.
\end{theorem}
\begin{proof}
Maintain a visited table on $1,\ldots,t$. Starting at the least unvisited
representative, apply inverse multiplication until it returns to that
representative. A permutation orbit cannot merge into a previously completed
orbit; otherwise applying inverse steps would put its starting point in that
orbit as well. Mark each encountered vertex, including those beyond the
window. The resulting disjoint orbits exhaust all representatives, each
visited once. Reject an orbit if any vertex exceeds $M$. Theorem~\ref{thm:classification}
proves that the retained orbits, and only those, reconstruct positive raw
cycles. At $M=0$ there are no positive cycles without any graph traversal.
Each retained H block has length $\ell+1\le L$. Disjointness gives a total
of at most $t$ retained H vertices and hence at most $tL$ raw instructions.
\end{proof}

The state-visit bound is not a total bit-complexity bound.
Repeated gcd removal takes at most $\lfloor\log_2 s\rfloor$ rounds, because
the product of the removed factors is $s$ and each is at least two.
For $c=2^a,D=2$, there are exactly $a$ rounds while $t=1$ remains fixed.
Modular inversion is another preprocessing operation. Big-integer division,
multiplication, gcd, and output copying have parameter-dependent bit costs.
Explicit raw expansion also has its own output cost: $c=1$ gives a cycle
with $t=h=1$ but raw period $D$, arbitrarily large.

The \href{../go/AdditiveCycles.go}{exact executable} stores parameters,
representatives, and periods with arbitrary precision. Enumeration uses a
machine-sized budget supplied by the caller. A capacity refusal or
callback interruption is an unfinished search, not a complete zero-cycle
answer. An empty window is a separate algebraic zero-cycle certificate.
Certificates record the canonical modular orbit, its H period, effective
modulus, and grouped raw period. Verification checks every edge, first-access
index, absence of an early root return, and every claimed field without raw
replay. Callback certificates are borrowed; retained copies own their
integer storage.

\section{Boundary calculations and formal proofs}

For $c=10,D=6$, the decomposition is $s=2,t=5$ and the window is $M=2$.
The representatives 1 and 2 give raw cycles $2\to12\to2$ and
$4\to14\to24\to4$, with periods 2 and 3. Representative 3 is excluded.
Indeed scaling the legal $c=5,D=6$ cycle starting at 3 by 2 would start at 6,
which immediately divides to 1 rather than following that scaled word.

For $c=5,D=2$, the modular orbit $(1,3,4,2)$ has indices $(1,1,0,0)$.
It expands to $1\to6\to3\to8\to4\to2\to1$, with raw period 6.
The grouped formula gives $4+10/5=6$; termwise integer division would fail.
For $c=10,D=4$, $M=1$ but the modular successor of 1 is 4, so no orbit
survives merely because its initial vertex lies in the window.
For $c=8,D=6$, $N=3$ but $M=0$: a nonempty core need not contain a cycle.
For $c=5,D=6$, representative 5 is residue zero and gives order one and
raw period 6; the positive zero-class representative is retained.

The \href{../lean/AdditiveCycles.lean}{Lean development} follows the same
steps, with the following principal declarations.
\begin{center}
\begin{tabular}{ll}
\toprule
Mathematical step & Lean declaration\\
\midrule
Factor extraction & \texttt{AdditiveOffsetDecomposition}\\
Coprime-factor maximality & \texttt{AdditiveCoprimeFactorMaximal}\\
Closed endpoint support & \texttt{AdditiveClosedPathSupported}\\
Closed endpoint core & \texttt{AdditiveClosedPathCore}\\
Exact scaled window & \texttt{AdditiveScaledCoreIff}\\
Guarded edge recovery & \texttt{AdditiveModularEdgeLegal}\\
Exact raw segmentation & \texttt{AdditiveEndingCorePath}\\
Return phase transport & \texttt{IterateReturnAlongCycle}\\
Raw-to-modular primitive map & \texttt{AdditiveRawPrimitiveToModular}\\
Primitive lift and clock & \texttt{AdditiveModularPrimitiveRaw}\\
Effective order & \texttt{AdditivePrimitiveOrder}\\
Finite length bound & \texttt{AdditivePrimitiveLengthBound}\\
Permutation edges & \texttt{AdditiveResiduePermutationEdge}\\
Positive representatives & \texttt{AdditivePositiveResidueUnique}\\
\bottomrule
\end{tabular}
\end{center}
The declaration \texttt{AdditiveRawPrimitiveClassification} combines the
forward data in one modular witness carrying its clock, order, and bound.
The cycle correspondence up to phase is assembled from these maps and the
endpoint recognition argument above; no separate phase-quotient API is used.

The test suite compares 25,506 cycles across all 4,512 pairs
$1\le c\le96$, $2\le D\le48$ with independently extracted raw cycles in
$1,\ldots,cD$. The executable enumerates and verifies the same catalogue
using its own certificate checker. A separate certificate has raw period
$2^{200}$ and is checked
algebraically, not by traversing that many instructions. These finite checks
validate the executable within their stated scope; the theorems above and
the formal proofs establish the universal positive cycle classification.
Neither signed parameters nor zero offset are asserted here.

\bibliographystyle{amsalpha}
\bibliography{AdditiveCycles}
\end{document}
